\documentclass[11pt,a4paper]{article}
\usepackage[utf8]{inputenc}
\usepackage[dutch]{babel}
\usepackage{csquotes}
\usepackage{amssymb}
\usepackage{enumitem}
\usepackage{epigraph}
\usepackage{setspace}
\usepackage{parskip}
\usepackage{xcolor}
\usepackage{hyperref}

\setlength{\parindent}{0pt}
\setlength{\parskip}{1em}
\renewcommand{\epigraphflush}{flushright}
\renewcommand{\epigraphsize}{\small}
\setlength{\epigraphwidth}{0.8\textwidth}

\title{\textbf{Het verkeerde hoekje} \\ \large \textit{Over begrensde modellen, en waarom daar een ethiek uit volgt}}
\author{}
\date{}

\begin{document}

\maketitle

\section*{I. Waar zit het heilige?}

Het Nederlandse woord \textit{god} is een oud woord, en het brengt ons in eerste instantie niet naar een ding. Het is van oorsprong een voltooid deelwoord. Wat het benoemde was iets wat gedaan werd: het aangeroepene, of datgene waaraan geplengd werd. In het Germaans was het bovendien onzijdig --- een \textit{wat}, geen \textit{wie}. Pas met de kerstening werd het mannelijk.

Ik gebruik dit niet als bewijs. De wortel is betwist, en zelfs als de etymologie klopt, volgt daar geen enkele filosofische stelling uit. Wat mij eraan bevalt is smaller: het woord wees oorspronkelijk naar een handeling, en pas later naar een substantie. Misschien is dat ook een betere plaats om het heilige te zoeken.

Ik ben opgegroeid in een gereformeerd gezin, in een dorp waar de kerk in de jaren tachtig nog een scheidslijn onder de bevolking was. Wat ik daar heb gezien, en wat ik nog steeds met waardering beschrijf, is een gemeenschap die goed was in samen leren. Mensen die elkaar bevroegen, hielpen, corrigeerden, en die generaties lang een gedeelde manier van omgaan met elkaar in stand hielden. Het is niet waar dat daar niet werd nagedacht. Er werd voortdurend nagedacht.

Op één punt niet.

Dat is mijn diagnose, en ik denk dat hij preciezer is dan de gebruikelijke kritiek. Het probleem met een onaantastbare hoeksteen is niet in de eerste plaats dat hij misschien fout is. Het probleem is dat hij niet meer kan deelnemen aan het proces waarmee je ontdekt of hij fout is. Alles mocht worden herzien behalve dat ene, en daarmee ontstond precies in het midden een blinde vlek. Op den duur wordt die vlek wat mensen niet meer kunnen geloven. Ik denk dat de kerk daaraan leegloopt: niet aan haar antwoorden, maar aan de uitzondering die zij op haar eigen werkwijze maakte.

De heiligheid zat in het verkeerde hoekje.

\section*{II. Waarom is er hoe dan ook een procedure nodig?}

Ik denk dat we in een universum leben waarin nieuwe dingen ontstaan. Niet alleen nieuwe gebeurtenissen binnen bestaande soorten, maar nieuwe soorten. Uit chemie komt leven, uit leven komt zenuwstelsel, uit zenuwstelsel komt taal, uit taal komt wetenschap, en uit wetenschap komen dingen waarvoor we nu nog geen woord hebben. Dat is wat ik met emergentie bedoel.

Het zwakke argument hier zou zijn: de wereld verandert, dus je kunt hem niet voorspellen. Dat klopt niet. Planeetbanen veranderen ook en zijn uitstekend voorspelbaar. Verandering op zichzelf is geen probleem.

Het sterke argument is een ander. Als er werkelijk nieuwe soorten dingen ontstaan, dan loopt mijn model niet alleen achter op de feiten --- het mist de vakken om die feiten in onder te brengen. Een model is een indeling van de werkelijkheid in categorieën. Wat buiten al mijn categorieën valt, kan ik niet als ontbrekend registreren, want ik heb geen plek waar het zou moeten staan.

Dat maakt het nog niet kenbare tot één domein met twee eigenschappen tegelijk. Elk afzonderlijk stuk ervan is in principe leerbaar --- daarom is wetenschap zinvol en niet ijdel. En het raakt nooit op, want er komt nieuw bij. Leren blijft zinvol \textit{én} eindigt nooit.

Wat daarmee permanent is, is geen ding. Het is de afstand tussen model en werkelijkheid.

\section*{III. Waarom hebben we elkaar nodig?}

Dit is de dragende stap, en het is de stap die ik het langst verkeerd heb gehad.

Mijn eerste formulering was: een systeem dat alleen interne feedback kent, kan zichzelf niet ijken. Dat is te grof, en het is aantoonbaar onwaar. Ik werk dagelijks met modellen die zichzelf uitstekend corrigeren. Kruisvalidatie, toetsing op achtergehouden data, het vergelijken van voorspelling met waarneming --- dat is allemaal een model dat zijn eigen fouten vindt zonder dat er iemand van buiten aan te pas komt.

De juiste claim is subtieler. Een systeem kan fouten vinden binnen zijn eigen correctieprocedure. Het kan niet vanuit zichzelf vaststellen of die procedure zelf toereikend is. Interne feedback vindt fouten; externe feedback kan de voorwaarden ter discussie stellen waaronder je fouten zoekt.

En daar sluit het vorige deel op aan. Interne correctie kan alleen fouten vinden in categorieën die het model al heeft. Emergentie levert nieuwe categorieën. De grens van zelfcorrectie is dus precies de grens van mijn categorieënset --- en dat is een grens die ik van binnenuit niet kan zien liggen, omdat ik hem alleen zou kunnen waarnemen met de categorieën die hem vormen.

Ook hier moet ik nog één keer preciezer zijn, want \enquote{extern} betekent niet vanzelf \enquote{een ander mens}. Het betekent eerst: informatie die niet door dezelfde representatie is voortgebracht. Een meetinstrument is in die zin extern. Maar een meting levert mij een waarde binnen een vak dat ik al had. Een ander mens kan mij een vak aandragen dat ik niet had --- een ander model, andere waarnemingen, een andere indeling. Dat is niet de enige externe bron. Het is de enige externe bron die het soort correctie kan leveren waar emergentie om vraagt.

Dat is de reden waarom anderen in dit betoog geen morele toevoeging zijn maar een structurele voorwaarde.

\subsection*{De keukentafel}

Ik heb dit mechanisme niet bedacht en daarna herkend. Ik heb het eerst meegemaakt en pas veel later kunnen formuleren.

In mijn huwelijk liep het gesprek vast wanneer ik een positie innam die niet paste bij hoe de ander de situatie zag. Dat is niet bijzonder; dat gebeurt overal. Wat wel bijzonder was: wanneer ik voorstelde om te bespreken \textit{hóé} we dit soort gesprekken voerden, liep dat ook vast.

Er zijn twee niveaus. Op het eerste zijn we het oneens. Op het tweede zijn we het oneens over hoe we met onenigheid omgaan. Zolang het tweede niveau openstaat, is het eerste oplosbaar --- dan kun je de procedure bijstellen en het opnieuw proberen. Sluit ook het tweede niveau, dan is er geen hoger intern niveau meer waarop het conflict kan worden opgelost.

Dat is dezelfde structuur als hierboven, maar dan geleefd. En het is waar het van een interessant probleem een dwingend probleem werd. Een gedachte-experiment heeft een uitgang. Gedeeld ouderschap heeft er geen.

\subsection*{Blijven}

Er is nog een voorwaarde, en die is niet epistemisch maar lichamelijk.

Uit de biologie kende ik de vlucht- en vechtreactie: een acute drang om weg te zijn. Ik had liefde nooit eerder ervaren als een reactie van dezelfde orde --- een acute reden om te blijven. Maar dat is wat het is. Het brein kent niet alleen een perfecte illusie van eigen gelijk, het kent ook reflexen die ons sturen. Het kan ons laten rennen. Het kan ons ook laten stoppen en teruglopen.

Correctie vereist aanwezigheid. Je kunt niet samen leren met iemand bij wie je bent weggelopen. Blijven is daarmee geen deugd naast het leren, maar de voorwaarde die samen leren mogelijk maakt.

Dit is de enige schakel in de hele keten die ik heb gevoeld voordat ik hem afleidde. Ik weet niet of dat een aanbeveling is of juist een reden tot achterdocht.

\subsection*{De test}

Elke stap hierboven gaat over modellen, niet over mensen. Er staat nergens iets in dat alleen voor biologische wezens geldt.

Dat is geen futuristische uitweiding maar de toets op mijn eigen redenering. Als deze ethiek alleen zou gelden voor menselijke subjecten, dan was de afleiding kennelijk niet structureel en heb ik ergens ongemerkt iets menselijks binnengesmokkeld. Geldt zij voor elk systeem dat de werkelijkheid modelleert --- mens, buitenaards, kunstmatig --- dan is het geen cultuur maar structuur.

\section*{IV. Wat gebeurt er als we dat verhinderen?}

In de traditie waarin ik ben opgegroeid was er een hel. De sanctie lag buiten het systeem: iemand anders hield de rekening bij, en het loonde pas later.

Hier ligt de sanctie binnen het systeem. Wie correctie afsluit, sluit een route naar nieuwe kennis af. De eerste structurele consequentie daarvan is vertraging. Een systeem dat zichzelf minder goed kan corrigeren, verliest zijn vermogen om zich aan te passen aan een werkelijkheid die blijft veranderen. De werkelijkheid hoeft daarvoor niets terug te doen. Er is geen boekhouder nodig.

Ik zeg bewust niet dat vertraging de enige consequentie is. Gesloten correctie leidt in de praktijk ook tot verkeerde beslissingen, tot schade, soms tot instorting. Maar de eerste en meest structurele sanctie is verlies van snelheid, en dat is precies genoeg --- het maakt de ethiek minder moralistisch, omdat er niets dreigt behalve wat er hoe dan ook gebeurt.

En de vertraging is wederkerig. Wie zwijgt waar hij zou moeten corrigeren, laat niet alleen de ander in een fout zitten. Hij houdt ook zijn eigen fouten uit het zicht, want hij verkleint het kanaal waarlangs zijn eigen correctie binnenkomt. Dat is geen beleefdheidsregel. Het is de reden waarom dit een ethiek is en geen methodologie.

\subsection*{Hier maak ik een sprong}

Ik wil die sprong niet verstoppen, want hij bepaalt hoe serieus de rest genomen kan worden.

Uit begrensdheid volgt niet dat we elkaar moeten ontzien. Uit begrensdheid volgt alleen dat we elkaar nodig hebben om niet blind te blijven. Dat ik dat vervolgens een reden vind om elkaar niet te vernietigen, is een waardeoordeel. Het volgt niet uit de structuur; ik voeg het toe.

Maar het is een kleiner waardeoordeel dan het lijkt. Ik hoef niet te vinden dat kennis heilig is, of dat waarheid boven alles gaat, of dat de mens goed is. Ik hoef alleen te vinden dat het beter is om mogelijkheden tot correctie open te houden dan ze te sluiten.

Er bestaat een oudere zet om die stap alsnog dwingend te maken: wie beargumenteert dat hij anderen mag vernietigen, doet daarmee een waarheidsclaim, en een waarheidsclaim veronderstelt de procedure die hij op datzelfde moment schendt. Ik vind dat elegant. Ik weet ook dat het betwist is, en ik leun er daarom niet op. Ik noem het alleen om te laten zien hoe smal de sprong is die ik maak.

\section*{V. Wat betekent dat voor een religieuze erfenis?}

Ik heb religie lang gelezen als een verzameling stellingen waar een praktijk uit voortkomt. Dat is niet hoe het werkte in de traditie waarin ik opgroeide. Daar was navolging niet het gevolg van het geloof --- het \textit{wás} het geloof. Je bent geen christen omdat je een propositie accepteert. Je bent het in het steeds opnieuw doen als Jezus.

Ik stel voor om dat als een aanknopingspunt te lezen in plaats van als een ontmaskering. Ik beweer niet dat ik weet wat religie werkelijk betekende, of dat ik haar ware kern heb blootgelegd. Ik zeg alleen dat ik, toen ik keek naar wat ik van huis uit had meegekregen, geen stelsel van beweringen vond maar een stelsel van werkwoorden. Herhalen, navolgen, samenkomen, elkaar bevragen, doorgeven.

Die werkwoorden heb ik niet uitgevonden. Ik heb ze geërfd. Wat ik heb gedaan is het object vervangen: niet een persoon om na te volgen, maar het leren zelf. De vorm bleef intact. Daarom is de hoop die ik eruit haal ook niet nieuw geconstrueerd --- hij is meeverhuisd.

Wat wel anders is: bij navolging van een voorbeeld zit het onbekende in de weg. Je weet waar je heen gaat en niet hoe ver je komt. Hier zit het onbekende ook in de bestemming. Emergentie levert categorieën die nu nog niet bestaan, dus wat \enquote{beter} over honderd jaar betekent, kan ik vandaag niet formuleren. Dat is niet minder hoop. Het is opener hoop, en juist daardoor overdraagbaar naar mensen die niets van mijn traditie delen.

Zonder garantiegever is elkaar het enige waar je op terug kunt vallen. Dat is geen sentimentele toevoeging aan het slot; het is de conclusie van deel III. En het heeft een eigenschap die ik pas laat heb opgemerkt: deze hoop is decentraal. Hij ligt niet bij één instantie, dus er is ook geen instantie die hem kan intrekken. Geen kerkscheuring, geen teleurstelling en geen enkel mens kan hem van je afnemen. Wat hem in stand houdt is niet vertrouwen in een uitkomst, maar het blijven doen van de werkwoorden.

Iets soortgelijks geldt voor de andere zekerheden waar mensen op leunen. Geld valt niet weg wanneer er geen garantiegever meer is --- geld is zelf precies zo'n gedeelde afspraak. Wat wegvalt is de gegarandeerde zekerheid: iets wat draagt zonder dat iemand het draagt. Wat overblijft is onderhouden zekerheid. Instituties, afspraken, relaties, vertrouwen --- allemaal even hard als het werk dat erin gaat.

\section*{VI. Wat betekent dat voor mij?}

Er is een verschil tussen begrijpen dat je een model bent en zo kunnen leven. Ik heb lang naar een beeld gezocht en heb er nu een.

Ik was een rots. Zo zag ik mezelf: iets wat vaststaat, waar anderen zich aan kunnen vasthouden, dat hetzelfde blijft. Ik ben een rollende steen geworden. Zwaar genoeg om te dragen, maar in beweging, en niet te voorspellen waar hij ligt.

Wat het beeld voor mij goed doet, is dat het gewicht behoudt. Traagheid is ook een vorm van zekerheid. Dat alle aannames in een model kunnen veranderen, betekent niet dat er niets is. En het brengt de rest van dit betoog terug in het lichaam: op een oppervlak dat rolt kun je alleen staan als er iets buiten je is om je aan vast te houden. Dat is dezelfde afleiding als in deel III, maar dan als evenwicht in plaats van als argument.

De belangrijkste wending kwam later. De rots is nooit een rots geweest. Ik stond altijd al op iets wat rolde; ik dacht alleen dat het vaststond. Wat er is veranderd is niet mijn positie maar mijn zicht erop. Ik heb dus geen stabiliteit ingeleverd. Ik heb een illusie van stabiliteit ingeleverd.

Daarom is \enquote{instabiel} ook niet het goede woord. Instabiel is iets wat omvalt. Wat ik beschrijf is voortdurend corrigeren om overeind te blijven --- en dat is wat lopen is, en fietsen. Een fiets is in natuurkundige zin instabiel en volkomen betrouwbaar zolang je trapt. Correctie is daar geen defect. Het is de manier waarop het evenwicht überhaupt tot stand komt.

En dan de zin waar het voor mij op uitkomt. Je staat niet op de steen. Je staat op het feit dat je weet dat hij rolt.

Van binnenuit voel ik dezelfde zekerheid als ieder ander, dus daar heb ik niets aan. Wat ik wél heb is kennis over hoe het instrument werkt, en die komt van buiten mijn hoofd: uit de biologie, uit het gedachte-experiment, uit anderen. Dat is geen zekerheid die je bezit. Het is een zekerheid die je onderhoudt.

Bob Dylan heeft een lied over een rollende steen, en daar is het verlies. Iemand die alles kwijtraakt wat als grond diende. Ik zou dat niet willen wegpoetsen, want dat verlies is echt. Maar de landing is anders. Bij hem blijft er niemand over. Hier is het juist het moment waarop anderen ophouden optioneel te zijn.

\section*{VII. En de balans dan?}

Hier houdt het essay op te helpen, en dat moet ik erbij zeggen.

In het echte leven is de vraag niet of correctie noodzakelijk is. De vraag is hoeveel zekerheid je vasthoudt en hoeveel correctie je toelaat, vandaag, in dit gesprek, met deze persoon. Daar geef ik geen formule voor. Ik heb er ook geen.

Dat is geen gat in het betoog. Het is de conclusie ervan. Bestond de formule, dan bestond er een autoriteit die hem kent, en dan viel de hele afleiding weg. Er is niemand die het weet, en dat is precies waarom er een procedure nodig is in plaats van een antwoord.

En dat maakt de diagnose uit deel I scherper dan ik hem eerst had. Een onaantastbare hoeksteen is nu juist het antwoord op deze spanning. Iemand legt de balans één keer vast, en daarna hoeft niemand hem meer per situatie te zoeken. Dat is een enorme verlichting, en het is aantrekkelijk om exact dezelfde reden waarom het misgaat. Wat ik hier aanbied is geen betere balans. Het is de vaststelling dat hij dagelijks opnieuw moet worden gevonden, en dat niemand dat werk van je overneemt. Dat is zwaarder, niet lichter.

Nog iets, en het is het punt dat ik zelf het langst over het hoofd heb gezien. De balans ligt niet binnen één persoon. Van binnenuit kun je niet vaststellen of je te vast staat of te veel meebeweegt --- beide voelen gerechtvaardigd, want het gevoel van gerechtvaardigd zijn is nu net wat niet betrouwbaar is. De balans wordt zichtbaar in de omgang met anderen. Zij zien wanneer ik hard ben geworden. Ik zie het bij hen.

Daarmee is evenwicht geen persoonlijke deugd maar opnieuw een relatie, en dat is dezelfde verschuiving als in deel I. En het verklaart wat er aan die keukentafel gebeurde. Het liep niet vast omdat iemand zijn balans verkeerd had afgesteld. Het liep vast omdat het kanaal waarlangs die balans wordt bijgesteld dicht zat. Zonder dat kanaal wordt iedereen vanzelf harder, want er is niets meer dat het bijstelt.

Eén waarschuwing aan mezelf. \enquote{Zekerheid in het proces} mag niet de nieuwe onaantastbare hoeksteen worden. Ook de procedure is herzienbaar; ook zij kan tekortschieten. Wat ik heb is bescheidener en daardoor steviger: van alle manieren die ik ken om met een begrensd model te leven, is dit de enige die haar eigen fouten kan vinden.

\section*{VIII. Kan ik mijn eigen theorie vertrouwen?}

Zeggen dat je het mis kunt hebben is te goedkoop. Vrijwel elke moderne tekst zegt dat, en generieke bescheidenheid is precies zo immuun voor correctie als een hoeksteen. Als dit essay iets waard is, moet het concreet aanwijzen waar het kan falen.

\begin{enumerate}[label=\textbf{\arabic*}.]
    \item Ik neem aan dat het beter kán worden. Emergentie garandeert dat niet. Zij levert nieuwe categorieën, en daaronder ook nieuwe manieren om het te verpesten. Wat wel volgt is dat leren de enige route is die beter mogelijk \textit{máákt}. Dat is minder dan wat ik zou willen zeggen.
    \item De stap van \textit{nodig hebben} naar \textit{mogen} is normatief. Ik heb hem in deel IV zichtbaar gemaakt en niet gedicht.
    \item De etymologie in deel I draagt niets. De wortel is betwist en ik gebruik haar als opening, niet als fundament. Wie hem eruit haalt, houdt het argument over.
    \item Dit essay weerlegt God niet en probeert dat ook niet. Wie zeker weet dat God niet bestaat, zit in dezelfde kamer met hetzelfde instrument als wie zeker weet dat hij wel bestaat. Wat hier wordt afgewezen is niet het geloof maar het hoekje: dat één ding immuun wordt gemaakt voor gesprek. Mijn ouders en ik denken nog steeds verschillend over God, en dat is geen breuklijn meer.
    \item De procedure zelf kan ontoereikend blijken, en ik kan dat van binnenuit niet vaststellen --- dat is letterlijk wat deel III beweert.
    \item Mijn herlezing van religieuze navolging in deel V is een voorstel, geen bewijs. Ik weet niet wat religie voor anderen is.
    \item \textbf{Zeven, en die weegt het zwaarst.} Het is mogelijk dat alles hierboven klopt en niets oplevert. Dat het meta-niveau bij de ander gesloten is, is geen uitzonderingsgeval --- het is de situatie waarin mensen het vaakst radeloos zijn. Dit betoog beschrijft die val correct en biedt geen uitweg wanneer de ander niet meedoet. Ik heb dat aan een keukentafel meegemaakt en ik heb er nog steeds geen antwoord op. Wie dit essay leest als een gereedschap voor zo'n situatie, zal het zien tekortschieten.
\end{enumerate}

En dan is er nog iets waar ik zelf niet omheen kan. Dit bouwwerk voelt van binnenuit sluitend. Alles haakt in elkaar, elke stap volgt uit de vorige, en dat is een prettig gevoel. Het is ook precies het gevoel waar het gedachte-experiment waar dit alles mee begon voor waarschuwt. Een systeem dat overal past voelt van binnenuit identiek aan een systeem dat klopt. Ik kan dat niet oplossen door beter na te denken, want denken is het instrument dat in het geding is.

Wat er dan overblijft is niet een houding maar een handeling. Ik moet dit voorleggen aan mensen met andere categorieën dan de mijne. Aan iemand die gelooft en die deel I ziet uitvoeren op zijn eigen hoeksteen. Aan iemand voor wie \enquote{we hebben elkaar nodig om elkaar te corrigeren} geen abstractie is maar een positie waarin hij ooit is uitgebuit. Aan mensen die iets tegenover dit betoog zetten in plaats van eraan mee te bouwen. Dat zijn de lezers die vakken hebben die ik niet heb.

Dus dit is geen slotwoord maar een verzoek, en het is het verzoek dat het hele essay uitvoert in plaats van beschrijft. Ik vraag niet alleen om kritiek op mijn conclusies. Ik vraag om kritiek op het mechanisme waarmee ik tot conclusies kom.

Want het heilige is dan uiteindelijk niet de hoeksteen, en ook niet de procedure. Het heilige is dat de weg terug naar correctie niet wordt afgesneden.

Daar zit mijn zekerheid, en daar zit mijn hoop. Niet in de uitkomst --- die ken ik niet en die kan ik niet kennen. Maar in het blijven doen van de werkwoorden.

Samen leren is samenleven.

\end{document}